\section{Annexe E -- Protocole V90}
Le protocole V90 unifie la géométrie dynamique, le temps propre local, la vitesse d'information et le redshift observé dans un cadre simple à comparer à $\Lambda$CDM.

\subsection{Objectif}
Établir la relation complète entre la géométrie dynamique, le temps propre local, la vitesse d'information, la vitesse effective et le redshift observé, afin de relier les anomalies cosmologiques: tension de Hubble, zone critique autour de $z\approx0.46$, supernovas et régimes compacts sans singularité.

\subsection{Hypothèses de base}
\begin{enumerate}
    \item La géométrie n'est pas infinie et reste régulière même dans les zones extrêmes.
    \item Le temps propre n'est pas universel: chaque région possède un rythme d'écoulement différent.
    \item La vitesse d'information dépend du temps propre local.
    \item La vitesse effective observée dépend de la géométrie et du temps propre.
    \item Le redshift contient une composante géométrique et une composante dynamique du photon.
\end{enumerate}

\subsection{Définition des quantités}
On note $\tau(z)$ le temps propre local, $d(z)$ la distance géométrique effective, $v_{\mathrm{info}}(z)$ la vitesse d'information, $v_{\mathrm{eff}}(z)$ la vitesse effective observée, $g(z)$ le facteur géométrique et $f_{\mathrm{photon}}(z)$ la contribution dynamique du photon.

\subsection{Déroulé calculé}
On déroule maintenant la chaîne de calcul utilisée dans le test V90.

\paragraph{Étape 1: fonction de transition}
On introduit une transition lisse autour de la zone critique:
\[
S(z)=\frac{1}{1+\exp\!\left(-\frac{z-z_c}{w}\right)}
\]
avec $z_c=0.46$ et $w=0.16$.

\paragraph{Étape 2: géométrie dynamique}
Le facteur géométrique est pris sous la forme
\[
g(z)=1+\alpha_g S(z),
\qquad \alpha_g=0.18.
\]
Ainsi, pour $z=0.46$, on a $S(0.46)=0.5$ et donc
\[
g(0.46)=1+0.18\times 0.5=1.09.
\]

\paragraph{Étape 3: temps propre local}
Le temps propre local suit directement
\[
\tau(z)=\tau_0 g(z),
\qquad \tau_0=1.
\]
Donc, à $z=0.46$:
\[
\tau(0.46)=1\times 1.09=1.09.
\]

\paragraph{Étape 4: distance géométrique standard}
La distance comobile standard est prise sous la forme
\[
d_{\mathrm{std}}(z)=c\int_0^z \frac{dz'}{H(z')}
\]
avec
\[
H(z)=H_0\sqrt{\Omega_r(1+z)^4+\Omega_m(1+z)^3+\Omega_\Lambda},
\]
et les paramètres du test $H_0=67.4$, $\Omega_m=0.315$, $\Omega_r=9\times10^{-5}$, $\Omega_\Lambda=0.685$.
Pour $z=0.46$, le calcul numérique donne
\[
d_{\mathrm{std}}(0.46)=0.27904193604302047.
\]

\paragraph{Étape 5: distance effective}
On introduit une compression locale de la distance:
\[
d(z)=\frac{d_{\mathrm{std}}(z)}{1+z}\,\frac{1}{1+\beta S(z)},
\qquad \beta=2.5.
\]
À $z=0.46$, cela donne
\[
d(0.46)=\frac{0.27904193604302047}{1.46}\,\frac{1}{1+2.5\times 0.5}
=0.08494427276804276.
\]

\paragraph{Étape 6: vitesse d'information}
La vitesse d'information locale suit alors
\[
v_{\mathrm{info}}(z)=\frac{d(z)}{\tau(z)}.
\]
Donc, à $z=0.46$:
\[
v_{\mathrm{info}}(0.46)=\frac{0.08494427276804276}{1.09}=0.07793052547526857.
\]

\paragraph{Étape 7: facteur photonique}
Pour imposer la relation de redshift testée dans le script, on définit
\[
f_{\mathrm{photon}}(z)=\frac{1+z}{g(z)}.
\]
Alors, à $z=0.46$:
\[
f_{\mathrm{photon}}(0.46)=\frac{1.46}{1.09}=1.3394495412844036.
\]

\paragraph{Étape 8: vitesse effective observée}
La vitesse effective observée est
\[
v_{\mathrm{eff}}(z)=v_{\mathrm{info}}(z)\,f_{\mathrm{photon}}(z).
\]
Donc, à $z=0.46$:
\[
v_{\mathrm{eff}}(0.46)=0.07793052547526857\times 1.3394495412844036
=0.10438400659990102.
\]

\paragraph{Étape 9: redshift total}
Le redshift total est
\[
1+z=g(z)f_{\mathrm{photon}}(z).
\]
À $z=0.46$:
\[
g(0.46)f_{\mathrm{photon}}(0.46)=1.09\times 1.3394495412844036=1.46,
\]
donc
\[
1+z=1.46 \quad\Rightarrow\quad z=0.46.
\]

\paragraph{Étape 10: cohérence numérique}
Le résidu de redshift vaut
\[
\Delta z = (g f_{\mathrm{photon}}-1)-z.
\]
Pour $z=0.46$, le test retourne
\[
\Delta z = -5.551115123125783\times 10^{-17},
\]
ce qui est numériquement nul.

\subsection{Résumé des valeurs testées}
Le calcul est stabilisé sur la grille $z=\{0,0.23,0.46,0.69,1,2,3\}$ avec les valeurs centrales suivantes:
\begin{center}
\begin{tabular}{lrrrr}
\hline
$z$ & $g(z)$ & $\tau(z)$ & $v_{\mathrm{info}}(z)$ & $v_{\mathrm{eff}}(z)$ \\
\hline
0.23 & 1.0345479015605026 & 1.0345479015605026 & 0.07883840064020031 & 0.0937329558555731 \\
0.46 & 1.09 & 1.09 & 0.07793052547526857 & 0.10438400659990102 \\
0.69 & 1.1454520984394974 & 1.1454520984394974 & 0.06717653676800872 & 0.09911226082050893 \\
\hline
\end{tabular}
\end{center}

On voit bien que le point critique $z\approx 0.46$ marque un ralentissement du temps propre, une baisse de la vitesse d'information et une augmentation du facteur photonique, tout en gardant l'identité de redshift exacte dans le test.

\subsection{Relations fondamentales}
\begin{align}
    \tau(z) &= \tau_0\,g(z),\\
    v_{\mathrm{info}}(z) &= \frac{d(z)}{\tau(z)},\\
    v_{\mathrm{eff}}(z) &= v_{\mathrm{info}}(z)\,f_{\mathrm{photon}}(z),\\
    1+z &= g(z)\,f_{\mathrm{photon}}(z).
\end{align}

En combinant ces relations, on obtient l'équation finale du modèle:
\[
1+z = \left[\frac{d(z)}{\tau_0\,g(z)}\right] f_{\mathrm{photon}}(z).
\]

\subsection{Zone critique}
Autour de $z\approx0.46$, le facteur géométrique varie rapidement, le temps propre ralentit, la vitesse d'information diminue et la vitesse effective augmente. Dans cette zone, le redshift devient plus fort que la prédiction standard.

\subsection{Conséquences physiques}
\begin{itemize}
    \item La vitesse d'information n'est pas constante.
    \item Le redshift n'est pas purement géométrique.
    \item La tension de Hubble devient une conséquence du couple $\tau(z)+g(z)$.
    \item Les supernovas super-lumineuses peuvent être relues comme une conversion gamma vers visible dans une géométrie extrême.
    \item Les trous noirs n'ont pas besoin de singularité si le temps propre diffère et si la dissipation géométrique reste finie.
\end{itemize}

\subsection{Suite du programme}
Les prochaines étapes sont de tracer $v_{\mathrm{info}}(z)$, définir explicitement $g(z)$, calibrer $f_{\mathrm{photon}}(z)$, puis comparer V90 à $\Lambda$CDM.